\chapter{The finite-volume scheme of the general rate model}
\label{app:grm}

Section~\ref{sec:grm} introduced the general rate model (GRM) and stated, without
proof, that its finite-volume discretisation is \emph{mass-conservative by
construction}. This appendix sets out the numerical method in full: the two
meshes and the discrete unknowns, the discrete bulk and particle balances, the
film coupling that ties them together, the fully implicit time step with its
Picard linearisation of the competitive isotherm, the assembly of the single
coupled linear system, and the mass-balance accounting used to verify the solver.
The implementation is \code{models/chromatography/grm.py}, which builds its
spatial operators with the \code{PyFVTool} finite-volume library and assembles the
coupled system through the helper module \code{models/chromatography/coupling.py}.

Throughout, $i=1,\dots,n_c$ indexes the protein components, $\eps$ and $\eps_p$ are
the bed and particle porosities, $u$ the interstitial velocity, $D_{\mathrm{ax}}$
the axial dispersion, $R_p$ the bead radius, $k_{f,i}$ the film coefficient, and
$D_{p,i}$ the pore diffusivity, exactly as in Section~\ref{sec:grm}. The continuous
problem is Eqs.~\eqref{eq:grm-bulk}--\eqref{eq:grm-bc}.

%==============================================================================
\section{Why finite volumes}
\label{app:grm-why}

The finite-volume method discretises the \emph{integral} form of a conservation
law: it divides the domain into non-overlapping control volumes (cells) and, for
each cell, balances the rate of change of the stored quantity against the fluxes
through the cell faces. Because a flux leaving one cell is, identically, the flux
entering its neighbour, whatever leaves one control volume is gained by the next,
and mass can be neither created nor destroyed by the discretisation---only
transported across faces or stored within cells. This local-telescoping property
is what makes the method the natural choice for a model whose entire reason for
existing (Section~\ref{sec:grm}) is to conserve a strongly retained, easily-lost
band. The method-of-lines engine of Section~\ref{sec:engine}, by contrast,
discretises the differential form and leaves conservation to the accuracy of the
time integrator---which is precisely what fails in the stiff regime.

%==============================================================================
\section{Two meshes and the discrete unknowns}
\label{app:grm-mesh}

The GRM lives on two coupled grids:

\begin{itemize}
  \item an \textbf{axial} Cartesian grid (\code{Grid1D}) of $N_z$ cells of width
  $\Delta z = L/N_z$ spanning the column, carrying the bulk concentrations
  $c_i$; and
  \item for \emph{each} axial cell $j$, a \textbf{radial} spherical grid
  (\code{SphericalGrid1D}) of $N_r$ shells spanning a representative bead,
  carrying the pore concentrations $c_{p,i}$.
\end{itemize}

Each \code{PyFVTool} cell variable carries two \emph{ghost cells} (one beyond each
boundary) that encode the boundary conditions algebraically, so a field on $N$
interior cells is stored as a vector of length $N+2$. The discrete unknowns are
therefore the bulk values $c_{i,a}$ ($a=1,\dots,N_z$ interior, plus two ghosts per
component) and the pore values $c_{p,i,j,k}$ ($k=1,\dots,N_r$ interior, plus two
ghosts per $(i,j)$ bead). Counting interior cells, the protein system carries
$n_c\,(N_z + N_z N_r)$ unknowns; the unretained modulator $m$ adds a further $N_z$
on the axial grid and is advanced separately.

A geometric subtlety, flagged in Section~\ref{sec:grm}, is settled here once. On
the spherical grid \code{PyFVTool} stores a midpoint cell volume
$4\pi r_c^2\,\Delta r$, but its diffusion operator conserves against the
\emph{true} shell volume
%
\begin{equation}
V_k = \tfrac{4}{3}\pi\bigl(r_{k+1/2}^3 - r_{k-1/2}^3\bigr),
\label{eq:app-shell}
\end{equation}
%
where $r_{k\pm1/2}$ are the shell face radii. The GRM uses $V_k$ of
Eq.~\eqref{eq:app-shell} for \emph{all} mass accounting and for the film
normalisation below, which removes the only inconsistency that would otherwise
spoil exact conservation.

%==============================================================================
\section{The discrete bulk balance}
\label{app:grm-bulk}

Integrating the bulk equation~\eqref{eq:grm-bulk} over axial cell $a$ and dividing
by the cell volume gives, for each interior cell,
%
\begin{equation}
\dv{c_{i,a}}{t}
  = -\,\frac{(F^{\mathrm{conv}}_{a+1/2} - F^{\mathrm{conv}}_{a-1/2})}{\Delta z}
    + \frac{(F^{\mathrm{disp}}_{a+1/2} - F^{\mathrm{disp}}_{a-1/2})}{\Delta z}
    - \underbrace{\frac{1-\eps}{\eps}\,a_v\,k_{\mathrm{ov},i}\,
      \bigl(c_{i,a} - c_{p,i,a,N_r}\bigr)}_{\text{film sink}},
\label{eq:app-bulkdisc}
\end{equation}
%
with $a_v = 3(1-\eps)/R_p$ the external specific area. The two transport terms are
built by \code{PyFVTool} operators:

\begin{itemize}
  \item \textbf{Convection} uses \emph{first-order upwinding}
  (\code{convectionUpwindTerm}): the face concentration is taken from the
  upstream cell, $F^{\mathrm{conv}}_{a+1/2}=u\,c_{i,a}$ for $u>0$. Upwinding adds a
  small numerical dispersion but is monotone---it cannot manufacture the negative
  concentrations that destroy the method-of-lines engine.
  \item \textbf{Dispersion} uses a central two-point gradient
  (\code{diffusionTerm}): $F^{\mathrm{disp}}_{a+1/2}=D_{\mathrm{ax}}\,
  (c_{i,a+1}-c_{i,a})/\Delta z$.
\end{itemize}

The inlet is a Dirichlet condition $c_i=c_{\mathrm{feed},i}(t)$ imposed through
the left ghost cell; the outlet is a zero-gradient (Danckwerts) condition. The
film sink---the last term of Eq.~\eqref{eq:app-bulkdisc}---is deferred to
Section~\ref{app:grm-film}, because it is shared, identically, with the particle
balance.

%==============================================================================
\section{The discrete particle balance}
\label{app:grm-particle}

Within bead $j$, integrating the particle equation~\eqref{eq:grm-part} over shell
$k$ gives a balance between the change of \emph{stored} solute and the diffusive
fluxes through the inner and outer shell faces,
%
\begin{equation}
V_k\,\dv{s_{i,j,k}}{t}
  = \eps_p D_{p,i}\Bigl(A_{k+1/2}\,\tfrac{c_{p,i,j,k+1}-c_{p,i,j,k}}{\Delta r}
                       - A_{k-1/2}\,\tfrac{c_{p,i,j,k}-c_{p,i,j,k-1}}{\Delta r}\Bigr),
\label{eq:app-partdisc}
\end{equation}
%
where $A_{k\pm1/2}=4\pi r_{k\pm1/2}^2$ are the shell face areas and
%
\begin{equation}
s_{i} \;=\; \eps_p\,c_{p,i} + (1-\eps_p)\,\qstar_i(\bm c_p, m, \mathrm{pH})
\label{eq:app-storage}
\end{equation}
%
is the \emph{stored} solute per unit particle volume---mobile pore liquid plus
adsorbed phase---evaluated with the \emph{same} isotherm $\qstar$ as everywhere
else in the package. The spherical radial Laplacian of
Eq.~\eqref{eq:app-partdisc} is exactly \code{PyFVTool}'s spherical
\code{diffusionTerm}. The centre face $k=\tfrac12$ has zero area, which enforces
the symmetry condition $\partial c_{p,i}/\partial r|_{0}=0$ automatically; the
outer face $k=N_r+\tfrac12$ carries the film flux of the next section.

%==============================================================================
\section{The film coupling as one shared discrete flux}
\label{app:grm-film}

The film boundary condition~\eqref{eq:grm-bc} sets the flux that crosses the bead
surface. The scheme combines the two outer resistances---the film proper and the
diffusion across the outermost half-shell---in series, giving an effective
coefficient
%
\begin{equation}
k_{\mathrm{ov},i}
  = \left(\frac{1}{k_{f,i}} + \frac{\Delta r/2}{\eps_p D_{p,i}}\right)^{\!-1},
\label{eq:app-kov}
\end{equation}
%
and defines a single discrete surface flux
$\Phi_{i,j}=k_{\mathrm{ov},i}\bigl(c_{i,a(j)}-c_{p,i,j,N_r}\bigr)$ between the bulk
cell $a(j)$ and the outer pore cell $N_r$ of its bead. The decisive design choice
is that the \emph{same} number $\Phi_{i,j}$ enters two places with opposite sign:

\begin{itemize}
  \item as a \textbf{sink} in the bulk cell, the film term of
  Eq.~\eqref{eq:app-bulkdisc} with $k_{\mathrm{ov},i}$ in place of $k_{f,i}$; and
  \item as a \textbf{source} on the bead surface cell, the outer-face flux of
  Eq.~\eqref{eq:app-partdisc}, normalised by the true outer shell volume
  $V_{N_r}$ of Eq.~\eqref{eq:app-shell}.
\end{itemize}

Because what the bulk loses through $\Phi_{i,j}$ the bead gains exactly---the same
discrete quantity, not two separately discretised approximations of it---the
bulk$\leftrightarrow$bead exchange conserves mass to machine precision regardless
of mesh resolution. In the assembled matrix this is four entries per
$(i,j)$ pair: $+$ on the bulk diagonal and $-$ on the bulk--surface off-diagonal,
$+$ on the surface diagonal and $-$ on the surface--bulk off-diagonal. The helper
\code{coupling.Assembler.couple} stamps exactly this quartet, so the sign
bookkeeping lives in one tested place rather than being rewritten per model.

%==============================================================================
\section{Implicit time stepping and the Picard linearisation}
\label{app:grm-time}

The scheme is advanced by \emph{backward Euler} (fully implicit): all spatial
terms are evaluated at the new time level $n+1$, which is unconditionally stable
and lets the stiff film and adsorption terms be taken in a single robust step. The
difficulty is that the stored solute~\eqref{eq:app-storage} depends nonlinearly,
and \emph{competitively}, on the new pore concentrations through
$\qstar_i(\bm c_p)$: every component's loading depends on every component's
concentration.

The model handles this in \emph{storage form} with a Picard (fixed-point)
linearisation. The exact implicit particle balance for shell $k$ reads
%
\begin{equation}
\frac{s_i(\bm c_p^{\,n+1}, m^{n+1}) - s_i(\bm c_p^{\,n}, m^{n})}{\Delta t}
  = \mathcal{D}_i\bigl[\bm c_p^{\,n+1}\bigr],
\label{eq:app-implicit}
\end{equation}
%
with $\mathcal{D}_i$ the discrete diffusion operator of
Eq.~\eqref{eq:app-partdisc}. Linearising the new-time storage about the current
Picard iterate $\bm c_p^{*}$ through the \emph{capacity Jacobian}
%
\begin{equation}
C_{i\ell} \;=\; \pdv{s_i}{c_{p,\ell}}
  \;=\; \eps_p\,\delta_{i\ell} + (1-\eps_p)\,\pdv{\qstar_i}{c_{p,\ell}},
\label{eq:app-capacity}
\end{equation}
%
and rearranging so the exact storage stays on the right-hand side gives the linear
update solved at each iteration,
%
\begin{equation}
\sum_\ell \frac{C_{i\ell}}{\Delta t}\,c_{p,\ell}^{\,n+1} - \mathcal{D}_i\bigl[\bm c_p^{\,n+1}\bigr]
  = \frac{1}{\Delta t}\Bigl(\sum_\ell C_{i\ell}\,c_{p,\ell}^{*}
        - s_i(\bm c_p^{*}, m^{n+1}) + s_i(\bm c_p^{\,n}, m^{n})\Bigr).
\label{eq:app-picard}
\end{equation}
%
The competitive cross-terms $C_{i\ell}$ ($i\neq\ell$) are obtained by a cheap
finite-difference perturbation of the isotherm and couple the components within
each shell. The construction has two properties worth stating explicitly, both
exploited deliberately:

\begin{keybox}{Why the storage form conserves mass exactly}
At Picard convergence $\bm c_p^{\,n+1}=\bm c_p^{*}$, so the linearised terms
$\sum_\ell C_{i\ell}c_{p,\ell}$ cancel between the two sides of
Eq.~\eqref{eq:app-picard} and what remains is precisely the \emph{exact} implicit
balance~\eqref{eq:app-implicit}: the change in true stored mass equals the net
diffusive flux. No separate $\dd q/\dd t$ source is introduced, so no
discretisation of the adsorption rate can leak mass. Moreover the \emph{old}
storage $s_i(\bm c_p^{\,n},m^{n})$ is evaluated at the \emph{old} modulator
$m^{n}$, so a moving salt gradient is differenced exactly rather than frozen---the
single device that lets one term capture both the nonlinear isotherm and the
changing gradient. For a linear isotherm $C$ is constant and one solve suffices;
otherwise the iteration is repeated to a relative tolerance (typically a handful
of steps).
\end{keybox}

The implementation of Eq.~\eqref{eq:app-picard} is the helper
\code{coupling.conservative\_storage}; the diagonal accumulation it reduces to for
the bulk field and the modulator is \code{PyFVTool}'s \code{transientTerm}.

%==============================================================================
\section{Assembling and solving the coupled system}
\label{app:grm-assembly}

\code{PyFVTool}'s own solvers target a single field on a single mesh. The GRM
instead couples $n_c$ bulk fields and $n_c N_z$ particle fields across two mesh
types into \emph{one} sparse system, so the discrete operators of the preceding
sections must be placed into a common global matrix. This index bookkeeping is the
job of \code{coupling.CoupledSystem}: each field is registered and handed a
contiguous block of the global unknown vector, and an \code{Assembler} scatters
every local operator block (convection, dispersion, diffusion, boundary terms,
the transient diagonal) to its block while \code{couple} stamps the inter-field
film quartet of Section~\ref{app:grm-film}. The geometric (constant) part of the
matrix---transport, pore diffusion, boundary structure, and the film
couplings---is assembled \emph{once}; only the feed-dependent inlet condition and
the state-dependent capacity blocks~\eqref{eq:app-capacity} are rebuilt each step
and added to it.

The resulting system $\bm M\,\bm x = \bm b$ is a single large, sparse,
\emph{non-symmetric} linear system (non-symmetric because of upwind convection and
the asymmetric film coupling), solved by a sparse \textsc{lu} factorisation
(\code{scipy.sparse.linalg.spsolve}). A direct solve is appropriate here: the
system is one-dimensional in structure and moderately sized, and the
factorisation is robust for the stiff, ill-conditioned blocks the film and
adsorption terms produce. \code{PyFVTool}'s \code{solveMatrixPDE} is \emph{not}
used, as it reshapes a solution onto a single mesh and so does not apply across the
coupled bulk-and-particle meshes.

\begin{keybox}{One time step, end to end}
\begin{enumerate}\setlength\itemsep{2pt}
  \item Advance the unretained modulator $m$ on the axial mesh (a small linear
  solve, accumulation by \code{transientTerm}).
  \item Freeze the old storage $s_i(\bm c_p^{\,n}, m^{n})$ at the old modulator.
  \item Picard loop until converged: evaluate the capacity blocks $C_{i\ell}$ at
  the current iterate; assemble the dynamic terms (feed inlet, bulk accumulation,
  conservative storage) onto the cached static matrix; solve the coupled sparse
  system for $\bm x^{n+1}$.
  \item Accumulate the inlet and outlet boundary fluxes for the mass balance,
  reading the inlet dispersive gradient from \code{PyFVTool}'s \code{gradientTerm}.
\end{enumerate}
\end{keybox}

%==============================================================================
\section{Mass-balance accounting and verification}
\label{app:grm-massbalance}

The solver reports a single dimensionless figure of merit, the relative
mass-balance closure error
%
\begin{equation}
\mathcal{E}
  = \frac{\;m_{\mathrm{in}} - m_{\mathrm{out}} - \Delta m_{\mathrm{stored}}\;}
         {m_{\mathrm{in}}},
\label{eq:app-massbal}
\end{equation}
%
where $m_{\mathrm{in}}$ and $m_{\mathrm{out}}$ are the time-integrated inlet and
outlet fluxes (carried in the void at the interstitial velocity, the inlet
including its dispersive contribution) and $\Delta m_{\mathrm{stored}}$ is the
total accumulated solute---bulk liquid plus, weighted by the true shell volumes
$V_k$ of Eq.~\eqref{eq:app-shell} and the number of beads per axial cell, the
pore and adsorbed phases. Equation~\eqref{eq:app-massbal} is an independent
audit: it is computed from the boundary fluxes and the final state, not used in
the solve, so a small $\mathcal{E}$ genuinely certifies conservation.

In practice $\mathcal{E}$ sits at the $10^{-14}$--$10^{-11}$ level---machine
precision---across the single-component breakthrough, the strongly-bound
multi-component isocratic load, and the competitive gradient elution that
motivated the model, confirming the construction of
Sections~\ref{app:grm-film}--\ref{app:grm-time}. This is the quantitative content
of Figure~\ref{fig:grm-cons}: where the method-of-lines engine loses essentially
all of the injected mass, the finite-volume GRM closes its balance to the last
few digits of floating-point.

\begin{remark}
The cost is the price of fidelity. Each step solves a coupled implicit system in
$n_c(N_z + N_z N_r)$ unknowns, possibly several times per step through the Picard
loop. Caching the static matrix, assembling the dynamic terms by sparse triplets
rather than dense slicing, and using a direct factorisation keep this practical,
but the GRM remains the reference solver for mechanism and conservation studies
rather than the tool for rapid design-space sweeps---for which the lumped engine
of Section~\ref{sec:engine} is retained.
\end{remark}
